% Tóm tắt, Phân chia công việc. TV1 ghép, cả nhóm cùng sửa.
\section*{Tóm tắt}
\addcontentsline{toc}{section}{Tóm tắt}
\phutrach{TV1, viết sau cùng}
Một hệ thống quản lý vận đơn phải ghi liên tục, đồng thời tra cứu theo nhiều thuộc tính, nên bảng vận đơn thường mang nhiều chỉ mục phụ. MySQL InnoDB lưu mỗi chỉ mục phụ thành một cây B+Tree trên một máy; TiDB biến mỗi mục chỉ mục thành một cặp khóa--giá trị riêng, phân tán trên các node TiKV và nhân bản bằng Raft. Báo cáo này đo chi phí của chỉ mục phụ trên hai hệ ở ba mặt: dung lượng, độ trễ đọc và thông lượng ghi.

TiDB v8.5.5 (1 PD, 1 TiDB server, 3 TiKV) và MySQL 8.4.11 chạy luân phiên trong Docker trên cùng một máy, với cùng hồ sơ tài nguyên XS: mỗi hệ 4 CPU và khoảng 750 MiB bộ nhớ đệm dữ liệu. Dữ liệu vận đơn 1 triệu và 3 triệu dòng được sinh giống nhau trên ba máy từ cùng một giá trị seed. TiDB được đo với 3 bản sao mỗi Region (C3) và, để tách chi phí nhân bản, với 1 bản sao (C1). Mỗi phép đo lặp lại nhiều vòng; báo cáo dùng trung vị kèm khoảng dao động và không loại lượt nào vì kết quả trái giả thuyết.

Về dung lượng, mô hình dựng từ định dạng lưu trữ khớp số đo trong khoảng 1,2\% với MySQL và tới từng byte với TiDB, trừ chỉ mục UNIQUE; mỗi mục chỉ mục của TiDB lớn hơn MySQL 3,4 đến 4,8 lần, tính cho một bản sao, trước khi nén. Về đọc, ở 3 triệu dòng covering index giúp TiDB nhiều hơn khi ít dòng khớp (nhanh gấp 2,01 lần so với 1,08 lần trên MySQL khi $R = 10$), nhưng giúp MySQL nhiều hơn khi nhiều dòng khớp (5,91 so với 3,31 lần khi $R = 10.000$). Quét bảng bắt đầu nhanh hơn dùng chỉ mục ở khoảng 3,2\% số dòng trên TiDB và 6,9\% trên MySQL. Dùng mã vận đơn làm khóa chính làm chỉ mục phụ lớn thêm 27,4\% trên MySQL và 39,5\% trên TiDB. Về ghi, thêm 5 chỉ mục làm thông lượng thêm đơn giảm 18,9\% trên MySQL nhưng 82,7\% trên TiDB C3, và TiDB C1 không nhân bản vẫn giảm 79,4\%. Như vậy giả thuyết ``LSM-tree giúp TiDB giảm ít hơn'' bị bác bỏ, và nhân bản không phải nguyên nhân chính. Với covering index, cái giá khi ghi phụ thuộc vào việc cột thêm vào có bị cập nhật hay không: khi \code{status} bị cập nhật, thông lượng giảm 29,5\% trên MySQL và 61,6\% trên TiDB C3.

Các kết quả này có giới hạn rõ. Mọi node chạy trên một máy và một ổ đĩa, nên không có độ trễ mạng vật lý và chênh lệch giữa hai hệ không được quy cho mạng; khả năng chịu lỗi của C3 chỉ là chịu mất một tiến trình TiKV. Quy mô dừng ở 3 triệu dòng, con số 10 triệu dòng chỉ là ngoại suy từ mô hình; thí nghiệm ghi chỉ dùng một mức 64 luồng; và số tuyệt đối chỉ được so giữa hai hệ trên cùng một máy.

\clearpage
\section*{Phân chia công việc và mức độ đóng góp}
\renewcommand{\thetable}{P.\arabic{table}}
\addcontentsline{toc}{section}{Phân chia công việc và mức độ đóng góp}
\ghichu{Mỗi người sửa dòng của mình theo việc đã làm thật (đối chiếu lịch sử Git). Nhóm thống nhất tỉ lệ đóng góp trước khi nộp.}

\needspace{9\baselineskip}
{\bangfont
\begin{longtable}{|>{\raggedright\arraybackslash\bfseries}p{2.9cm}|>{\raggedright\arraybackslash}p{3.0cm}|>{\raggedright\arraybackslash}p{7.2cm}|>{\centering\arraybackslash}p{1.5cm}|}
\caption{Phân chia công việc và mức độ đóng góp của các thành viên}\label{tab:p1}\\
\hline
\dongdau \textbf{Thành viên} & \textbf{Vai trò} & \textbf{Công việc cụ thể} & \textbf{Tỉ lệ} \\
\hline
\endfirsthead
\multicolumn{4}{c}{\textbf{Bảng P.1.} Phân chia công việc và mức độ đóng góp (tiếp theo)}\\[2pt]
\hline
\dongdau \textbf{Thành viên} & \textbf{Vai trò} & \textbf{Công việc cụ thể} & \textbf{Tỉ lệ} \\
\hline
\endhead
\endfoot
\endlastfoot
\TenTVMot & Trưởng nhóm, máy M1 & Mô hình và số đo dung lượng (B.1); RQ1 ở 3 triệu dòng; kiến trúc TiDB; ghép báo cáo, slide & 33,3\% \\ \hline
\TenTVHai & Môi trường, máy M2 & Cấu hình TiDB và MySQL, công cụ đo dùng chung; RQ4; RQ3; code RQ1; cài đặt và toàn bộ demo A.3; B.3 & 33,3\% \\ \hline
\TenTVBa & Chi phí đọc, máy M3 & RQ1 ở 1 triệu dòng; RQ2; thực nghiệm sơ bộ B.0; B.2 & 33,3\% \\ \hline
\end{longtable}
}

\needspace{9\baselineskip}
{\bangfont
\begin{longtable}{|>{\raggedright\arraybackslash\bfseries}m{3.6cm}|>{\raggedright\arraybackslash}m{11.6cm}|}
\caption{Chia việc viết báo cáo}\label{tab:p2}\\
\hline
\dongdau \textbf{Thành viên} & \textbf{Phần phụ trách} \\
\hline
\endfirsthead
\multicolumn{2}{c}{\textbf{Bảng P.2.} Chia việc viết báo cáo (tiếp theo)}\\[2pt]
\hline
\dongdau \textbf{Thành viên} & \textbf{Phần phụ trách} \\
\hline
\endhead
\endfoot
\endlastfoot
\TenTVMot & Tóm tắt; Chương 1; Chương 3 (vấn đề, mục tiêu, phạm vi, cơ chế); Chương 4 (dữ liệu và dung lượng); Chương 7 (giả thuyết, khuyến nghị); Kết luận \\ \hline
\TenTVHai & Chương 2; Chương 3 (phương pháp, kịch bản); Chương 4 (RQ3); Chương 6 \\ \hline
\TenTVBa & Chương 4 (thực nghiệm sơ bộ); Chương 5; Chương 7 (hạn chế) \\ \hline
\end{longtable}
}
